\section{Anforderungsanalyse}
In diesem Kapitel werden die Anforderungen an den zu entwickelnden Prototypen beschrieben. 
Die Anforderungen ergeben sich aus der Zielsetzung, zu untersuchen, wie Attribute Stream-Based Access Control im Linux-Kernel prototypisch umgesetzt werden kann.
Sie beschreiben das Verhalten, das der Prototyp im Rahmen der Arbeit zeigen soll.

\subsection{Funktionale Anforderungen}

Funktionale Anforderungen beschreiben, welche fachlichen Funktionen der Prototyp bereitstellen muss. 
Sie sind im Regelfall direkt testbar, da geprüft werden kann, ob eine Funktion vorhanden ist und das erwartete Verhalten zeigt. 
Im Kontext dieser Arbeit betreffen sie insbesondere das Einlesen und Verarbeiten von Policies, die Zugriffskontrolle über LSM-Hooks sowie die nachträgliche Überwachung bestehender Zugriffe.

\begin{description}
    \item[FA-01 Policy-Einlesen:]
    Der Prototyp muss ASBAC-Policies aus einer definierten Quelle im Userspace einlesen können.
    
    \textit{Akzeptanzkriterium:} Eine im Policy-Verzeichnis abgelegte Policy wird erkannt und für die weitere Verarbeitung bereitgestellt.

    \item[FA-02 Policy-Verwaltung:]
    Der Prototyp muss Änderungen an Policies erkennen und aktualisierte Policies in die Zugriffskontrolle übernehmen können.
    
    \textit{Akzeptanzkriterium:} Wird eine Policy geändert, hinzugefügt oder entfernt, wird diese Änderung vom System verarbeitet.

    %\item[FA-03 Kommunikation zwischen Userspace und Kernelspace:]
    %Der Prototyp muss eine Schnittstelle bereitstellen, über die Informationen zwischen Userspace und Kernelspace ausgetauscht werden können.
    
    %\textit{Akzeptanzkriterium:} Der Userspace kann dem Kernel entscheidungsrelevante Informationen übermitteln.

    \item[FA-04 Zugriffskontrolle bei Dateiöffnungen:]
    Der Prototyp muss Dateiöffnungen über den Linux-Security-Module-Hook \texttt{file\_open} abfangen können.
    
    \textit{Akzeptanzkriterium:} Das Öffnen einer durch eine Policy erfassten Datei wird durch den LSM-Hook \texttt{file\_open} kontrolliert.

    \item[FA-05 Entscheidungsfindung anhand von Policies:]
    Der Prototyp muss Zugriffsentscheidungen auf Grundlage der geladenen Policies treffen können.
    
    \textit{Akzeptanzkriterium:} Ein Zugriff wird abhängig von einer Policy erlaubt oder verweigert.

    \item[FA-06 Unterstützung dynamischer Attribute:]
    Der Prototyp muss dynamische Attribute, beispielsweise Zeitinformationen oder externe Zustandsinformationen, in Zugriffsentscheidungen einbeziehen können.
    
    \textit{Akzeptanzkriterium:} Eine Änderung eines dynamischen Attributs kann zu einer veränderten Zugriffsentscheidung führen.

    \item[FA-07 Überwachung und Enforcement bestehender Dateizugriffe:]
    Der Prototyp muss bereits bestehende Dateizugriffe überwachen und unmittelbar beenden, wenn sie durch entscheidungsrelevante Policy- oder Attributänderungen unzulässig werden.
    Als bestehender Zugriff gilt im Rahmen des Prototyps ein geöffneter File Descriptor eines laufenden Prozesses auf eine durch eine Policy erfasste Datei.
    Der Userspace-PEP ermittelt den zugehörigen Prozess und File Descriptor, bewertet den Zugriff anhand der aktuellen Policy- und Attributlage und schließt den File Descriptor bei einer Verletzung direkt.
    Eine reine Ausgabe der Verletzung ohne Enforcement ist nicht vorgesehen.

    \textit{Akzeptanzkriterium:} Öffnet ein Prozess eine Datei zunächst zulässig und wird dieser Zugriff anschließend durch eine Policy- oder Attributänderung unzulässig, schließt der Userspace-PEP den betroffenen File Descriptor. Andere File Descriptors desselben Prozesses bleiben geöffnet, sofern sie nicht selbst eine Policy verletzen.

    \item[FA-09 Administrationsschnittstelle:]
    Der Prototyp soll eine einfache Administrationsschnittstelle bereitstellen, über die zentrale Informationen zum Zustand des Prototyps lesend eingesehen werden können.
    
    \textit{Akzeptanzkriterium:} Ein Administrator kann zentrale Informationen zum Zustand des Prototyps abrufen, ohne über die Administrationsschnittstelle Policies, Attribute oder Kernelzustand zu verändern.

    \item[FA-10 Gültige Policy-Generationen:]
    \hypertarget{anf:gueltige-policy-generationen}{Der Prototyp muss sicherstellen, dass nur erfolgreich erzeugte und validierte Generationen von Policies für Zugriffsentscheidungen verwendet werden.}
    Fehlerhafte, unvollständige oder nicht erfolgreich verarbeitete Policy-Generationen dürfen nicht als aktive Entscheidungsgrundlage gelten.
    
    \textit{Akzeptanzkriterium:} Eine neue Policy-Generation wird erst dann aktiviert, wenn sie vollständig verarbeitet und erfolgreich validiert wurde. Andernfalls bleibt die zuletzt gültige Policy-Generation aktiv.

    \item[FA-11 Eingabevalidierung:]
    Der Prototyp muss Eingaben aus dem Userspace validieren, bevor diese im Kernel für Zugriffsentscheidungen verwendet werden.
    
    \textit{Akzeptanzkriterium:} Ungültige oder manipulierte Eingaben werden erkannt und nicht als Entscheidungsgrundlage im Kernel verwendet.

    \item[FA-12 beliebige Attributnamen:]
    In Policies muss es für den Nutzer möglich sein, Attributnamen frei zu vergeben, ohne dass diese vorher im Prototyp fest definiert sein müssen.
    Zulässige Attributnamen dürfen nicht leer sein und bestehen aus ASCII-Buchstaben (\texttt{A-Z}, \texttt{a-z}), Ziffern (\texttt{0-9}), Unterstrich (\texttt{\_}) und Bindestrich (\texttt{-}).
    
    \textit{Akzeptanzkriterium:} Eine Policy mit einem gültigen, zuvor nicht bekannten Attributnamen, beispielsweise \texttt{subject.position == \textquotedbl{}engineer\textquotedbl{}}, wird erfolgreich verarbeitet. Eine Policy mit einem Attributnamen außerhalb des zulässigen Zeichensatzes wird abgelehnt.

\end{description}

\subsection{Operative Anforderungen}

Operative Anforderungen beschreiben, wie gut der Prototyp seine Funktionen im Betrieb erfüllen soll. 
Sie beziehen sich auf Eigenschaften, die für Nutzung, Betrieb und Evaluation sichtbar sind, beispielsweise Stabilität, Beobachtbarkeit, Nachvollziehbarkeit, Performance und Reproduzierbarkeit.

\begin{description}
    \item[OA-01 Stabilität:]
    Der Prototyp darf die Stabilität des Linux-Kernels nicht unnötig gefährden.
    
    \textit{Akzeptanzkriterium:} Fehlerhafte oder ungültige Eingaben aus dem Userspace führen nicht zu einem Kernel-Absturz.

    \item[OA-02 Beobachtbarkeit und Nachvollziehbarkeit:]
    Der Prototyp soll zentrale Zustände und Entscheidungen so sichtbar machen, dass Betrieb und Evaluation nachvollziehen können, welche Policies, Attribute oder Zustände eine Zugriffsentscheidung beeinflusst haben.
    
    \textit{Akzeptanzkriterium:} Für zentrale Zugriffsentscheidungen ist erkennbar, welche Policy oder welcher Zustand die Entscheidung beeinflusst hat.

    \item[OA-03 Performance:]
    Die zusätzlichen Prüfungen durch den Prototyp sollen den kontrollierten Zugriff nicht unverhältnismäßig verlangsamen.
    
    \textit{Akzeptanzkriterium:} Die Auswirkungen auf die Ausführungszeit werden in der Evaluation betrachtet.

    \item[OA-04 Reproduzierbarkeit:]
    Die Evaluation des Prototyps soll unter dokumentierten Bedingungen wiederholbar sein.
    
    \textit{Akzeptanzkriterium:} Testumgebung, Policies und Testszenarien werden so beschrieben, dass die Ergebnisse nachvollzogen werden können.
\end{description}

\subsection{Entwicklungsbezogene Anforderungen}

Entwicklungsbezogene Anforderungen beschreiben Eigenschaften, die vor allem für die Entwicklung und spätere Erweiterung des Prototyps relevant sind. 
Sie sind für Nutzer nicht immer unmittelbar sichtbar, beeinflussen aber, wie verständlich, änderbar und erweiterbar das System bleibt.

\begin{description}
    \item[EA-01 Wartbarkeit:]
    Die Implementierung soll modular aufgebaut sein, sodass Userspace-Komponenten, Kernel-Komponenten und Kommunikationsschnittstelle getrennt betrachtet werden können.

    \item[EA-02 Erweiterbarkeit:]
    Der Prototyp soll so entworfen werden, dass weitere Attribute, Policies oder LSM-Hooks ergänzt werden können.

    \item[EA-03 Begrenzter Kernel-Anteil:]
    Komplexe Policy-Logik soll möglichst im Userspace verbleiben, um die Komplexität im Kernel zu reduzieren.
\end{description}




\newpage

